\section{Chem Core}
    \subsection{Matter and Measurement}
    Matter are elements of existence with mass and volume, that is it is made of stuff and takes up space. We use tools to make measurements and from that do science
        \subsubsection{Scientific Method}
        The scientific method is both a process and never ending cycle. It's basic structure is as follows;
        \begin{enumerate}
            \item Ask a Question
            \item Research Similar Problems
            \item Create a Hypothesis—an educated guess on what you expect the answer is
            \item Devise and Conduct and Experiment
            \item Analyze the data
            \item Make a conclusion
            \item Find your answer, or a new question
        \end{enumerate}
        The scientific method can begin at any step and creates an endless loop of discovery. It is a method of thinking based on logic and reason, asking questions and looking for answers through experimentation. The conclusion you come to is based on data and facts that can be proven
        \subsubsection{Error}
        Reality is messy and because of that there will be errors, no matter what happens perfection will never occur. However, in order to do science, they must be acknowledge and taken into the context of the results. In an experiment, you need to identify the three elements of an error: The error itself (what affected the data), the variable(s) affected (what numbers were affected, measured or calculated), and the magnitude of the error (how much did the results of the experiment change). There are three magnitudes of error, major, minor, and insignificant. Major errors changed the conclusion of the experiment. Instead of pressure being inversely correlated to volume, an error caused it to be found that it is proportional. Minor errors change the numbers, but not enough to change the conclusion. The ideal gas constant is found to be 5.23 J mol$^{-1}$ K$^{-1}$ instead of 8.31 J mol$^{-1}$ K$^{-1}$. Insignificant errors are errors that occurred but did not change the numbers themselves. Errors should be unavoidable, things like evaporation and temperature fluctuations, not mistakes such as incorrect measurements or improper technique. They also must be specific, not vague statements such as ``human error'' or ``inaccuracies'' but instead specific moments of error. 

        The way to measure the magnitude of the error of an experiment is with the percent error formula,
        \begin{equation}\label{eq:percentError}
            \text{Err}_\%=\frac{\text{Measurment}-\text{Accepted}}{\text{Accepted}}\times 100\%.
        \end{equation}
        Here `Measurement' is the value that you measured and `Accepted' is the value that many scientists have tested and confirmed to be accurate. If Err$_\%$ is negative, the measured value is lower than expected and if it is positive it is higher than expected.
        \subsubsection{Tools and Measurments}
        When using any tool to make a measurement, you will get a certain number of digits of accuracy, along with an additional digit that is obtained through estimation. Take the following thermometer,
        \begin{figure}[h]
            \centering
            \includegraphics[width=0.4\textwidth]{Screenshot 2026-05-31 at 10.17.01 PM.png}
        \end{figure}
        The tool tells us that the temperature is 21\textdegree C, but we must estimate the last digit since our measurement is between the two graduations. The final measurement should be around 21.6\textdegree C, with some wiggle room for the inaccuracies of estimation. 
        \subsubsection{Significant Digits}
        Significant Digits are the digits that we know are accurate, and their rules help us make accurate calculations and round appropriately while maintaining as much precision as possible. A digit is significant if, it is not zero, a zero between significant digits,  or a zero after a decimal point and after other significant digits. For AP Chemistry, you are correct if you are within once significant digit of accuracy. 
        \subsubsection{Extensive vs. Intensive Properties}
        A property is Extensive if it depends on the amount of a sample. Things like mass, volume, weight. A property is Intensive if it is independent on the amount of material present. Things like density or color. 
    \subsection{Particles and Reaction}
    Particles are tiny discrete pieces of matter, things like atoms, protons, neutrons, and electrons are all particles. We can measure the amount of these particles through their mass.

    \subsubsection{Atoms}
    Atoms are the particles that are core to chemistry. Atoms are the building blocks that make up everything, but are themselves made up of smaller particles called protons, neutrons and electrons. Atoms are composed of a dense nucleus at the center of the atom that contains the protons and neutrons and an electron cloud that makes up majority of the volume and houses the electrons.

    Electrons, denoted \ce{e-} have a negative charge and a low mass. Their low mass makes them susceptible to greater acceleration and therefore the move very fast in the electron cloud. Inside this cloud they exist at certain stable energy levels

    Protons, denoted \ce{p+}, are positively charged particles that live in the nucleus. They have a mass of 1 amu, much larger than the tiny electron. The number of protons in a nucleus determines the element through the atomic number, which equals the number of protons.

    Neutrons, denoted \ce{n^0}, are neutrally charged particles that live in the nucleus, they also have a mass of 1 amu. They help hold the nucleus together, acting as a sort of glue. There can be different numbers of neutrons for the same number of protons and the different configurations are called isotopes.

    Ions are atoms with a non-neutral charge due to an imbalance of electrons and neutrons. The charge is written as a superscript to the left of the atomic symbol. 

    Isotopes are the different variations of atoms with the same atomic number but different weights. They are written as the symbol-mass\# where the mass number is written in amu or atomic mass units. The average atomic mass of an element is the weighted average of the masses of its isotopes, based on their natural abundance. One graph that shows this is the mass spectra graph, which shows the abundance of each isotope based on the mass number.

    Atomic symbols can hold all of this information by including the mass number in the top left,the atomic number in the bottom left, and the charge in the top right.
    \[\ce{_6^12C^0}\]
    The above example represents a neutral Carbon-12 atom with 6 protons, 6 electrons, and 6 neutrons. In general the atomic symbol for any atom is given by
    \[_{\text{p}^+}^{\text{p}^+ + \text{n}^0}\text{X}^{\text{p}^+ - \text{e}^-}\]
    where the symbol X would be replaced to match the symbol associated with the atomic number. 

    \subsubsection{Bonding}
    Atoms form bonds to obtain a lower energy by filling their outer layer of electrons. This is because a full outer electron shell provides a more stable configuration for the atom. Nobel gases already have a full outer layer and that is why they are so unreactive. Atoms go about bonding in to different ways, Covalent Bonds and Ionic Bonds.

    In an Ionic bond one atom gives up some of its electrons to another atom forming two ions, a positive cation and a negative anion. In an ionic compound electrons are exchanged until both atoms have a full shell and the resulting compound has a ratio of ions resulting from their charges. The name of ionic compounds come from the name of the cation followed by the name of the anion changed so that it ends with -ide (oxide, fluoride, ect). Some cations have multiple charges that can be found in nature, when that happens we differentiate between them by putting roman numerals in the name. For example Iron has both 3+ and 2+ charges in nature, so we write \ce{Fe_2O_3} as Iron(iii) oxide since the iron ion \ce{Fe^3+} has a charge of 3+. There are also certain polyatomic ions that are made up of several atoms bound with covenant bonds such as \ce{NH4+} or \ce{NO3-} which have their own names. Below is a table containing the common polyatomic ions.
    \begin{table*}[h]
        \centering
        \begin{tabular}{c c c}
            Name & Chemical Formula & Charge \\
            \hline
            Ammonium & \ce{NH4+} & 1+ \\
            Acetate & \ce{C2H3O2-} & 1- \\
            Carbonate & \ce{CO3^2-} & 2- \\
            Chlorate & \ce{ClO3-} & 1- \\
            Chlorite & \ce{ClO2-} & 1- \\
            Chromate & \ce{CrO4^2-} & 2- \\
            Cyanide & \ce{CN-} & 1- \\
            Dichromate & \ce{Cr2O7^2-} & 2- \\
            Hydrogen Carbonate (Bicarbonate) & \ce{HCO3-} & 1- \\
            Hydrogen Sulfate & \ce{HSO4-} & 1- \\
            Hydride & \ce{H-} & 1- \\
            Hypochlorite & \ce{ClO-} & 1- \\
            Nitrate & \ce{NO3-} & 1- \\
            Nitrite & \ce{NO2-} & 1- \\
            Perchlorate & \ce{ClO4-} & 1- \\
            Permanganate & \ce{MnO4-} & 1- \\
            Peroxide & \ce{O2^2-} & 2- \\
            Phosphate & \ce{PO4^3-} & 3- \\
            Sulfate & \ce{SO4^2-} & 2- \\
            Sulfite & \ce{SO3^2-} & 2- \\
            Thiosulfate & \ce{S2O3^2-} & 2- \\
        \end{tabular}
        \caption{Polyatomic Ions Names and Chemical Formula}
    \end{table*}

    The charge taken by an atom is given by its oxidation number which can vary but does have a trend based on an elements position on the periodic table. Metals have positive oxidation numbers, transition metals typically have more than one oxidation number, non-metals and semi-metals have both positive and negative oxidation numbers, no element exists in a compound with an oxidation number greater than +8 and the most negative oxidation numbers equals the group number - 8. Nobel gases have no oxidation number as they do not form ions naturally. Below in \ref{table:OxidationNumbers} is a list of the first few elements and their possible oxidation numbers.
    
    Covalent bonds are formed when two atoms share electrons. The shared electrons allow both atoms to fill their outer shell and achieve a lower energy state. Covalent bonds can be single bonds, where one pair of electrons is shared, double bonds where two pairs of electrons are shared, and triple bonds where three pairs of electrons are shared. 
    \subsubsection{Advanced Compound Naming}
    The oxidation numbers of atoms also relate to the naming of polyatomic ions. The prefixes and suffixes for the polyatomic ion are based on the number of oxygen atoms and by extension the oxidation number of the central atom. The suffix -ate is used for the most common ion, -ite is used for the ion with one less oxygen, and hypo- and per- are used for ions with two less or two more oxygen atoms respectively. For example, the nitrate ion \ce{NO3-} has an oxidation number of +5 for nitrogen, so the nitrite ion \ce{NO2-} has an oxidation number of +3. The perchlorate ion \ce{ClO4-} has an oxidation number of +7, so the chlorate ion \ce{ClO3-} has an oxidation number of +5, the chlorite ion \ce{ClO2-} has an oxidation number of +3, and the hypochlorite ion \ce{ClO-} has an oxidation number of +1. The atoms that have polyatomic ions ending in the suffix -ate and 4 oxygen atoms form two squares called slivka squares, centered around niobium and arsenic and containing the 8 surrounding elements. All other atoms form -ate compounds with 3 oxygen atoms. The exception to this is Manganese which forms a -ate compound with 4 oxygen atoms, \ce{MnO4-}. 
    
    As hydrogen are added to these polyatomic ions, the name gains the prifix hydrogen or dihydrogen and the charge increases by 1 for each hydrogen added. For example, the carbonate ion \ce{CO3^2-} becomes the hydrogen carbonate ion \ce{HCO3-} with one hydrogen and the dihydrogen carbonate ion \ce{H2CO3} with two hydrogens. The sulfate ion \ce{SO4^2-} becomes the hydrogen sulfate ion \ce{HSO4-} with one hydrogen and the dihydrogen sulfate ion \ce{H2SO4} with two hydrogens. Another prefix that is commonly used is bi- which is used for ions with one hydrogen, such as the bicarbonate ion \ce{HCO3-} and the bisulfate ion \ce{HSO4-}.

    When naming an acid, the name of the acid depends on the name of the anion. If the anion ends in -ide, the acid name begins with hydro- and ends with -ic. For example, \ce{HCl} is hydrochloric acid because the anion is chloride. If the anion ends in -ite, the acid name ends with -ous. For example, \ce{HNO2} is nitrous acid because the anion is nitrite. If the anion ends in -ate, the acid name ends with -ic. For example, \ce{HNO3} is nitric acid because the anion is nitrate.

    Some ionic compounds also form complexes with water adding a third word to the name. Hydrates are ionic compounds that have a specific number of water molecules attached to them. The number of water molecules is indicated by a prefix and the word hydrate. For example, \ce{CuSO4*5H2O} is copper(II) sulfate pentahydrate because it has 5 water molecules attached to each formula unit of copper(II) sulfate.

    To name a covalent compound, we use the prefixes mono-, di-, tri-, tetra-, penta-, hexa-, hepta-, octa-, nona-, and deca- to indicate the number of atoms of each element in the compound. The first element is named using its full name, while the second element is named using its root and the suffix -ide. For example, \ce{CO2} is carbon dioxide because it has one carbon atom and two oxygen atoms. \ce{N2O4} is dinitrogen tetroxide because it has two nitrogen atoms and four oxygen atoms. The prefix mono- is typically omitted for the first element, so \ce{CO} is carbon monoxide, not monocarbon monoxide.

    In AP chemistry it is also expected that you know the names of the simple organic alkanes, which are hydrocarbons with only single bonds. The names of the alkanes are based on the number of carbon atoms in the molecule, with the suffix -ane. For example, \ce{CH4} is methane because it has one carbon atom, \ce{C2H6} is ethane because it has two carbon atoms, \ce{C3H8} is propane because it has three carbon atoms, and \ce{C4H10} is butane because it has four carbon atoms. The names of alkanes with more than four carbon atoms are based on Greek prefixes for the number of carbon atoms, such as pentane for five carbon atoms, hexane for six carbon atoms, and so on.

    \subsubsection{Periodic Trends}
    The periodic table is organized in a way that allows us to predict the properties of elements based on their position. The periodic trends are patterns in the properties of elements that occur as you move across a period or down a group in the periodic table. Some of the most important periodic trends include atomic radius, ionization energy, electronegativity, and electron affinity.

    Atomic radius is the size of the atom. The atomic radius increases down towards the left, with the largest atom being Francium. 

    Electronegativity is the ability of an atom to attract electrons in a chemical bond. Electronegativity increases up towards the right, with the most electronegative element being Fluorine.

    Ionization energy is the energy required to remove an electron from an atom. Ionization energy increases up towards the right, with the highest ionization energy being Helium.

    Electron affinity is the energy change that occurs when an atom gains an electron. Electron affinity increases up towards the right, with the highest electron affinity being Chlorine.

    In addition, we split the elements into metals, nonmetals, and metalloids. Metals are typically found on the left side of the periodic table and are characterized by their ability to conduct electricity, malleability, and ductility. Nonmetals are typically found on the right side of the periodic table and are characterized by their lack of conductivity, brittleness, and high ionization energies. Metalloids are elements that have properties of both metals and nonmetals and are typically found along the zig-zag line that separates metals and nonmetals on the periodic table.
    \subsection{Stoichiometry and the Mole}
        \subsubsection{Dimensional Analysis}
        Dimensional analysis is a method of converting between different units of measurement by using conversion factors. A conversion factor is a ratio of equivalent measurements that can be used to convert from one unit to another. To use dimensional analysis, you begin by writing the known value that you wish to convert, then multiply by a ratio with the unit you are canceling on opposite sides. For example to convert 5 kg to grams, you would write
\begin{equation*}
    5 \text{ kg} \times \frac{1000 \text{ g}}{1 \text{ kg}} = 5000 \text{ g}
\end{equation*}
        \subsubsection{Moles and Avogadro's Number}
        A mole is a unit of measurement that is used to count particles in chemistry. One mole is equal to 6.022 $\times$ 10$^{23}$ particles, which is called Avogadro's number. The mole is used to count atoms, molecules, ions, and other particles in chemistry. The molar mass of a substance is the mass of one mole of that substance, and it is typically expressed in grams per mole (g/mol). 
        \subsubsection{Standered Conditions}
        Many values in chemistry are measured at standered conditions, which are defined as a temperature of 0\textdegree C (273.15 K) and a pressure of 1 atm. At these conditions, one mole of an ideal gas occupies a volume of 22.4 liters. This is known as the molar volume of an ideal gas at standard temperature and pressure (STP).
        \subsubsection{Molarity, Density, Percent Composition}
        Concentration is a measure of the amount of solute dissolved in a given amount of solvent or solution. The most common unit of concentration is molarity (M), which is defined as the number of moles of solute per liter of solution. To find the molarity of a solution, you can use the formula:
        \begin{equation}\label{eq:molarity}
            M = \frac{n}{V}
        \end{equation}
        where $n$ is the number of moles of solute and $V$ is the volume of the solution in liters. This value is sometimes show in brackets. For example a 0.5 M solution of sodium chloride would be written as [NaCl] = 0.5 M.

        Density is a measure of how much mass is contained in a given volume. Density is not a measure of concentration because it uses the mass and volume of the entire solution, not just the solute. 

        The mole fraction can be used to convert between different concentrations. The mole fraction of a component in a solution is the ratio of the number of moles of that component to the total number of moles of all components in the solution. The mole fraction can be used to convert between molarity and mass percent, or between molarity and density.

        Parts per million (ppm) is a unit of concentration that is used to express very dilute concentrations. One ppm is equal to one part of solute per million parts of solution, which is equivalent to 1 mg of solute per liter of solution for aqueous solutions.

        Percent Composition is a way to express the concentration of a component in a solution as a percentage. It is calculated by dividing the mass of the component by the total mass of the solution and multiplying by 100\%.
        \subsubsection{Empirical and Molecular Formulas}
        The empirical formula for a compound is the ratio of atoms in a molecule. To find the empirical formula from percent composition or mass, divide each percent or mass by the corresponding atomic mass, then divide all by the smallest value to get the simplest whole number ratio. Rounding may be necessary to get whole numbers, and if the ratio is close to a common fraction such as 1.5, you can multiply all ratios by the denominator of that fraction to get whole numbers.
        To find the molecular formula from the empirical formula, you need to know the molar mass of the compound. Divide the molar mass by the empirical formula mass to get a whole number, then multiply the subscripts in the empirical formula by that whole number to get the molecular formula.
        \subsubsection{Hydrated Formulas}
        When trying to find the formula of a hydrate from mass or percent composition, you can use the same steps for finding the empirical formula, instead using the molar mass of water (18.015 g/mol) for the water component. The resulting formula will be in the form of \ce{X*YH2O} where X is the formula of the anhydrous compound and Y is the number of water molecules per formula unit of the anhydrous compound.
        \subsubsection{Dilutions}
        When diluting a solution, the number of moles remains constant while the concentration changes. The relationship between the initial and final concentrations and volumes can be expressed with the formula:
        \begin{equation}\label{eq:dilution}
            M_1V_1 = M_2V_2
        \end{equation}
        where $M_1$ and $V_1$ are the initial molarity and volume, and $M_2$ and $V_2$ are the final molarity and volume after dilution.